\documentclass[10pt]{article}
\usepackage[margin=0.72in]{geometry}
\usepackage[T1]{fontenc}
\usepackage{amsmath,amssymb,amsthm,mathtools,booktabs,array,microtype}
\usepackage{enumitem}
\usepackage{placeins}
\usepackage{xcolor}
\usepackage{hyperref}
\setlength{\parskip}{2.5pt}
\setlength{\parindent}{1.2em}
\setlist{nosep,leftmargin=1.6em}
\newtheorem{theorem}{Theorem}[section]
\newtheorem{proposition}[theorem]{Proposition}
\newtheorem{lemma}[theorem]{Lemma}
\newtheorem{corollary}[theorem]{Corollary}
\theoremstyle{definition}
\newtheorem{definition}[theorem]{Definition}
\newtheorem{remark}[theorem]{Remark}
\newcommand{\CI}[3]{[#1#2\mid #3]}
\newcommand{\A}{\mathcal A}
\newcommand{\Tfour}{T_4}
\newcommand{\supp}{\operatorname{supp}}
\newcommand{\AP}{\operatorname{ap}}
\newcommand{\code}[1]{\texttt{#1}}
\newcolumntype{P}[1]{>{\raggedright\arraybackslash}p{#1}}

\newcommand{\statusline}{Preprint, version 0.5}

\hypersetup{
  colorlinks=true,
  linkcolor=black,
  citecolor=black,
  urlcolor=blue,
  pdftitle={The antecedent threshold for five-variable Gaussian conditional-independence inference},
  pdfauthor={Maarten Linus Herrmann},
  pdfsubject={The antecedent threshold for five-variable regular-Gaussian Horn inference},
  pdfkeywords={Gaussian conditional independence; gaussoid; Horn implication; interval certification; Krawczyk operator}
}

\title{The antecedent threshold for five-variable Gaussian conditional-independence inference}
\author{Maarten Linus Herrmann\\
\small Erasmus School of Economics\\
\small Erasmus University Rotterdam\\
\small Rotterdam, the Netherlands\\
\small\href{https://orcid.org/0009-0005-8721-6588}{ORCID 0009-0005-8721-6588}}
\date{\statusline\\27 October 2026}

\usepackage{xurl}
\begin{document}
\maketitle
\vspace{-1.25em}

\begin{abstract}
We determine the least number of antecedents in a premise-minimal, valid elementary Horn implication with full five-label support that is not a semantic consequence of the complete marginal-and-conditional four-variable theory $T_4$.  The implications have a single elementary conclusion, and $T_4$ requires each of the ten four-variable marginal and conditional faces to be regular-Gaussian realizable.  For regular positive-definite Gaussians, the sharp threshold is $a_5=4$.  The upper bound uses the five-variable instance of \v{S}ime\v{c}ek's implication family and an 80-atom $T_4$ model satisfying its premises and violating its conclusion.  The new contributions are this nonlocality certificate, an exhaustive three-antecedent lower bound, and the resulting premise minimality.  The lower bound classifies premise triples under permutations and duality.  A positive-definite extension handles representatives with proper premise support; for full-support representatives, rational Krawczyk certificates give Gaussian counterexamples to every conclusion outside a computed sound closure.  An exact-arithmetic verifier checks the finite catalogues, orbit reductions, countermodel, interval certificates, and coverage.
\end{abstract}

\section{Introduction}

For a regular Gaussian vector, an elementary conditional-independence statement $X_i\perp X_j\mid X_K$ is equivalent to the vanishing of an almost-principal minor of its covariance matrix.  A Horn implication among such statements asks whether one minor must vanish whenever a specified collection of other minors vanishes throughout the positive-definite cone~\cite{Matus2005}.

The gaussoid axioms give universal local rules for these implications.  They characterize regular-Gaussian structures on three labels, but on four labels additional relations are needed.  Ln\v{e}ni\v{c}ka and Mat\'u\v{s} classified the four-variable case: of the 679 four-gaussoids, 629 are regular-Gaussian realizable, separated by five further implication schemes~\cite{LM2007}.  No finite collection of bounded-variable rules characterizes Gaussian conditional independence in all dimensions~\cite{Simecek2006,Sullivant2009}.

We measure the complexity of an implication by its number of distinct antecedents.  Boege proved that every regular-Gaussian inference form with at most two antecedents follows from the gaussoid axioms~\cite{Boege2022,BoegeCorrection2022}.  On five labels, we seek the least number of antecedents in a valid elementary Horn implication that is not a consequence of the complete four-variable theory.

Let $\Tfour$ denote the theory obtained by requiring all five marginal and all five conditional four-variable faces to belong to the 629-element regular-Gaussian catalogue.  Define $a_5$ as the least number of antecedents in a premise-minimal valid rule with full five-label support that is not a semantic consequence of $\Tfour$.  The main result is
\[
 \boxed{a_5=4}.
\]

The upper bound follows from a correlation proof of the five-variable instance of \v{S}ime\v{c}ek's implication.  An 80-bit $\Tfour$ model satisfies its four premises and violates its conclusion, proving that the rule is not four-local.  For the lower bound, all three-antecedent sets are reduced under permutations and duality.  A positive-definite extension proves the claim for proper-support representatives; rational interval certificates cover every conclusion outside the iterated $\Tfour$ closure for full-support representatives.

The result concerns elementary singleton-conclusion Horn implications for regular positive-definite Gaussians.  It does not classify all four-antecedent rules or all five-variable Gaussian CI structures.  Sections~\ref{sec:setup}--\ref{sec:lower} define the problem and prove the two bounds.  Section~\ref{sec:verification} summarizes the computational proof and methods, Section~\ref{sec:prior} places the threshold in the literature, and the appendices specify the countermodel, catalogue construction, and certificate data.

\section{Gaussian CI structures and the four-local theory}\label{sec:setup}

\subsection{Elementary atoms and almost-principal minors}

Let $N$ be a finite label set and let $X=(X_i)_{i\in N}$ be centered Gaussian with positive-definite covariance matrix $\Sigma$.  For distinct $i,j\in N$ and $K\subseteq N\setminus\{i,j\}$, write
\[
 [ij\mid K]\quad\text{for}\quad X_i\perp X_j\mid X_K,
\]
and abbreviate $[ij\mid\varnothing]$ by $[ij]$.  The elementary atom $[ij\mid K]$ is equivalent to
\begin{equation}\label{eq:ap}
 \AP_{ij\mid K}(\Sigma)
 :=\det \Sigma_{\{i\}\cup K,\,\{j\}\cup K}=0.
\end{equation}
The order of rows and columns in \eqref{eq:ap} affects only a fixed sign and never the zero condition.  The set of elementary atoms on $N$ is denoted by $\A_N$.  It has
\[
 |\A_N|=\binom{n}{2}2^{n-2};
\]
for $n=5$ this is $80$.

A \emph{regular-Gaussian CI structure} is the set of elementary atoms whose almost-principal minors vanish for some $\Sigma\succ0$.  Positive diagonal rescaling preserves all these zero patterns, so $\Sigma$ may be normalized to a correlation matrix $R=(r_{ij})$.  For a singleton conditioning set,
\begin{equation}\label{eq:partial}
 [ij\mid k]\quad\Longleftrightarrow\quad r_{ij}-r_{ik}r_{jk}=0.
\end{equation}

A \emph{gaussoid} is an elementary CI structure satisfying the four standard gaussoid axiom schemes~\cite{Matus2005,LM2007,Geometry2019}.  Every regular-Gaussian structure is a gaussoid, but the converse fails.

\subsection{Marginal and conditional four-faces}

Fix $N=\{1,2,3,4,5\}$.  For $x\in N$, deleting $x$ gives the \emph{marginal four-face}
\[
 \{[ij\mid K]:i,j\ne x,\ K\subseteq N\setminus\{i,j,x\}\}.
\]
The \emph{conditional four-face at $x$} is obtained by restricting to atoms that condition throughout on $x$ and then removing $x$ from the displayed conditioning sets:
\[
 [ij\mid K]\longleftrightarrow[ij\mid K\cup\{x\}],\qquad i,j\ne x.
\]
A regular Gaussian remains regular Gaussian after marginalization or conditioning, so every such face is a realizable four-variable structure.

The four-variable catalogue contains 679 gaussoids and exactly 629 regular-Gaussian structures~\cite{LM2007}; denote the latter set by $\Tfour^{(4)}$.  The computation enumerates all models of the elementary gaussoid axioms on the 24 four-variable atoms, then retains those satisfying the five additional Ln\v{e}ni\v{c}ka--Mat\'u\v{s} implication schemes.  The identification of the retained models with regular-Gaussian structures uses the classification theorem in~\cite{LM2007}.

\begin{definition}[The four-local theory $\Tfour$]\label{def:t4}
A five-label elementary CI structure is a model of $\Tfour$ if each of its ten marginal and conditional four-faces belongs to $\Tfour^{(4)}$.  A Horn implication is a semantic consequence of $\Tfour$ if it holds in every such 80-bit model.
\end{definition}

The computation also uses a constructive closure.  For each face, retain the members of $\Tfour^{(4)}$ containing the current face premises and adjoin every atom common to them; repeat over the ten faces to a fixed point.  Each added atom holds in every compatible four-variable model, so the procedure is sound for Definition~\ref{def:t4}.  Section~\ref{sec:lower} proves that every Gaussian-valid conclusion with three antecedents belongs to this closure.  A global 80-bit model of Definition~\ref{def:t4} satisfying the premises and omitting the conclusion refutes semantic consequence and hence every such derivation.

Every five-variable gaussoid-axiom instance is contained in one of these faces.  It has three distinguished labels and a background conditioning set $L$.  If $|L|\le1$, an unused label can be omitted; if $|L|=2$, conditioning the entire face on either member of $L$ contains the instance.  Thus every $\Tfour$ model is itself a gaussoid.

\subsection{Support, premise minimality, and the threshold}

For $a=[ij\mid K]$, set $\supp(a)=\{i,j\}\cup K$.  For a rule $P\Rightarrow c$, its support is the union of the supports of its premises and conclusion.

\begin{definition}[Genuine five-variable rule and $a_5$]\label{def:a5}
A Horn implication $P\Rightarrow c$ on five labels is \emph{genuinely five-variable beyond $\Tfour$} if:
\begin{enumerate}
 \item it is valid for every regular five-dimensional Gaussian vector;
 \item its support is all five labels; and
 \item it is not a semantic consequence of $\Tfour$.
\end{enumerate}
It is \emph{premise-minimal} if no proper subset of $P$ validly implies $c$.  Define $a_5$ as the minimum $|P|$ over premise-minimal genuinely five-variable rules.
\end{definition}

The identity covariance satisfies every elementary CI atom, so positive-premise sets are never vacuous.  Distinct antecedents are counted once; repeating an atom does not change $|P|$.

\section{The threshold theorem}\label{sec:main}

\begin{theorem}[Antecedent threshold]\label{thm:main}
Four antecedents are necessary and sufficient for a genuinely five-variable elementary regular-Gaussian CI Horn implication beyond $\Tfour$.  Equivalently,
\[
 \boxed{a_5=4}.
\]
\end{theorem}

The theorem follows from the next two propositions.

\begin{proposition}[Four-antecedent upper bound]\label{prop:upper}
The implication
\begin{equation}\label{eq:cycle5}
 [12\mid3]\wedge[13\mid4]\wedge[14\mid5]\wedge[15\mid2]
 \quad\Longrightarrow\quad[12]
\end{equation}
is regular-Gaussian valid, has support all five labels, and is not a consequence of $\Tfour$.
\end{proposition}

\begin{proposition}[Three-premise lower bound]\label{prop:lower}
Every regular-Gaussian-valid elementary Horn implication on five labels with at most three distinct antecedents is a semantic consequence of $\Tfour$.
\end{proposition}

Proposition~\ref{prop:upper} and Corollary~\ref{cor:minimal} give $a_5\le4$, including the premise-minimality clause of Definition~\ref{def:a5}.  Proposition~\ref{prop:lower} gives $a_5\ge4$.

\section{The four-antecedent upper bound}\label{sec:upper}

The implication family used here is due to \v{S}ime\v{c}ek~\cite{Simecek2006}; its five-variable member is \eqref{eq:cycle5}.  The following direct calculation proves the Gaussian-validity part of Proposition~\ref{prop:upper}.

Set
\[
 a=r_{12},\quad b=r_{13},\quad c=r_{14},\quad d=r_{15},
 \qquad p=r_{23},\quad q=r_{34},\quad u=r_{45},\quad s=r_{25}.
\]
By \eqref{eq:partial}, the four premises in \eqref{eq:cycle5} are
\[
 f_1=a-bp=0,\qquad f_2=b-cq=0,\qquad
 f_3=c-du=0,\qquad f_4=d-as=0.
\]
They satisfy the polynomial identity
\begin{equation}\label{eq:telescoping}
 f_1+p f_2+pq f_3+pqu f_4=a(1-pqus).
\end{equation}
Since every $2\times2$ principal correlation submatrix is positive definite,
\[
 |p|,|q|,|u|,|s|<1.
\]
The second factor on the right of \eqref{eq:telescoping} is therefore positive.  Under the four premises its left side vanishes, so $a=r_{12}=0$.  This proves the implication without division by a correlation.

The same telescoping calculation proves \v{S}ime\v{c}ek's general regular-Gaussian implication.  For distinct labels $0,1,\ldots,k$, with $k\ge3$,
\begin{equation}\label{eq:family}
 [01\mid2]\wedge[02\mid3]\wedge\cdots\wedge[0,k-1\mid k]
 \wedge[0k\mid1]\quad\Longrightarrow\quad[01].
\end{equation}
Write $a_j=r_{0j}$, $p_j=r_{j,j+1}$ for $1\le j<k$, and $p_k=r_{1k}$.  The premises give
\[
 a_j-a_{j+1}p_j=0\quad(1\le j<k),\qquad a_k-a_1p_k=0.
\]
After multiplying the $j$th equation by $p_1\cdots p_{j-1}$ and summing, one obtains
\[
 a_1\bigl(1-p_1p_2\cdots p_k\bigr)=0.
\]
Positive definiteness gives $|p_j|<1$ for every $j$, hence $a_1=0$.  The case $k=3$ is also rule (20) of Ln\v{e}ni\v{c}ka and Mat\'u\v{s}~\cite{LM2007}; the case $k=4$ gives Proposition~\ref{prop:upper}.

\section{Nonlocality and premise minimality}\label{sec:nonlocal}

To prove that the rule is not a consequence of the four-variable theory, we give a finite countermodel.

\begin{proposition}[$\Tfour$ nonlocality]\label{prop:nonlocal}
There exists an 80-bit model of all ten $\Tfour$ face constraints that contains the four premises of \eqref{eq:cycle5} and omits the atom $[12]$.  Hence \eqref{eq:cycle5} is not a semantic consequence of $\Tfour$.
\end{proposition}

\begin{proof}
The bitstring in Appendix~\ref{app:atoms} satisfies the following conditions, checked exhaustively.
\begin{enumerate}
 \item Bits $1,10,20,25$, corresponding to the atoms $[12\mid3]$, $[13\mid4]$, $[14\mid5]$, and $[15\mid2]$, are one; bit $0$, corresponding to the atom $[12]$, is zero.
 \item Every instance of the four elementary gaussoid axiom schemes holds for all ordered triples and background conditioning sets.
 \item For each omitted label, the marginal face in canonical 24-atom order belongs to the 629-element catalogue.
 \item For each label, the conditional face in the same order belongs to the catalogue.
\end{enumerate}
All ten masks belong to $\Tfour^{(4)}$.  The bitstring is therefore a global model of Definition~\ref{def:t4} and a countermodel to the implication.
\end{proof}

The countermodel is an abstract gaussoid with realizable four-faces.  Since it violates the valid implication~\eqref{eq:cycle5}, it has no five-variable Gaussian realization.

\begin{corollary}[Premise minimality]\label{cor:minimal}
The four antecedents in \eqref{eq:cycle5} are premise-minimal.
\end{corollary}

\begin{proof}
The full premise set does not force the atom $[12]$ in $\Tfour$.  By monotonicity, no three-premise subset forces it in $\Tfour$ either.  Proposition~\ref{prop:lower} implies that a Gaussian-valid deletion would be $\Tfour$-local, a contradiction.  Thus every deletion is Gaussian-invalid.  All four deletions canonicalize to the premise orbit represented by atom indices $(1,10,20)$, and certificate 272 covers the four mapped conclusion indices $0,8,16,24$.
\end{proof}

\section{The three-antecedent lower bound}\label{sec:lower}

For Proposition~\ref{prop:lower}, we classify the finitely many three-antecedent premise sets.  The cases with fewer antecedents follow from Boege's two-antecedental completeness theorem.  Proper-support representatives are handled analytically; full-support representatives require certified Gaussian counterexamples.

\subsection{Finite enumeration and symmetry reduction}

There are $\binom{80}{3}=82{,}160$ sets of three distinct atoms.  We quotient them by
\[
 G=S_5\times\langle *\rangle,
 \qquad [ij\mid K]^*=[ij\mid N\setminus(K\cup\{i,j\})].
\]
Label permutations preserve Gaussian validity and $\Tfour$.  For a positive-definite covariance matrix, inversion realizes duality on its CI structure.  This follows from Jacobi's complementary-minor identity
\[
 \det\bigl((\Sigma^{-1})_{I,J}\bigr)
 =\pm\frac{\det(\Sigma_{J^c,I^c})}{\det\Sigma},
\]
applied to the almost-principal minors.  Duality also exchanges marginal and conditional four-faces, while inversion preserves the 629-element four-variable catalogue.  Thus both Gaussian validity and semantic $\Tfour$ consequence are invariant under the 240 atom maps in $G$.

The enumeration selects one lexicographically canonical representative $P$ from each premise orbit.  For fixed $P$, conclusions are further quotiented by the stabilizer $G_P$.  If $P\Rightarrow c$ is invalid or $\Tfour$-local, applying any $g\in G$ gives the same alternative for $gP\Rightarrow gc$; no rule is lost by these two orbit reductions.  The constructive closure is also equivariant: $C(gP)=gC(P)$, since $G$ permutes the ten faces and preserves their catalogues.

Premise support is not invariant under duality, because a previously absent label can enter a conditioning set.  The terms \emph{proper support} and \emph{full support} in Table~\ref{tab:counts} therefore refer to the selected representative.  There are 105 orbits containing triples of both types.  Since Gaussian validity and $\Tfour$ consequence are orbit invariant, the conclusion proved for a representative holds for every member of its orbit.

\begin{table}[ht]
\centering
\caption{Three-antecedent classification counts.}
\label{tab:counts}
\begin{tabular}{P{5.5cm}rP{7.1cm}}
\toprule
Stage & Count & Meaning\\
\midrule
Labelled antecedent triples & 82,160 & All $\binom{80}{3}$ distinct three-atom sets\\
Premise orbits & 439 & Orbits under $S_5$ and duality\\
Canonical representatives with proper support & 104 & Representative premise support has at most four labels\\
Canonical representatives with full support & 335 & Representative premise support is all five labels\\
Mixed-support premise orbits & 105 & Orbit contains both proper- and full-support triples\\
Rule orbits in the 335-representative stratum & 21,357 & Orbits of conclusions $c\notin P$ under each representative stabilizer\\
Forced by iterated $\Tfour$ closure & 190 & Sound four-local derivations\\
Rule orbits outside the computed $\Tfour$ closure & 21,167 & Refuted by certified Gaussian witnesses\\
Rational interval certificates & 364 & One box for 306 representatives and two boxes for 29\\
\bottomrule
\end{tabular}
\end{table}

The 439 representatives split into 104 proper-support and 335 full-support cases.  Lemma~\ref{lem:extension} proves the claim for the proper-support cases.  For each full-support representative $P$, the verifier computes the fixed point $C(P)$ of the constructive $\Tfour$ closure from Section~\ref{sec:setup}.  Membership $c\in C(P)$ gives a sound four-local derivation of $P\Rightarrow c$.

It remains to rule out Gaussian validity for conclusions outside $C(P)$.  Quotienting the nontrivial conclusions $c\notin P$ by $G_P$ gives 21,357 rule orbits across the 335 representatives.  Of these, 190 lie in the computed closure and 21,167 do not.  Each of the latter is Gaussian-invalid, as shown by a certified witness.

For each $P$, a certificate records a set of conclusion atoms whose almost-principal minors are certified nonzero at its Gaussian root.  The union of these cover sets is checked to equal
\[
 \A_5\setminus C(P).
\]
This equality is checked on all 80 atoms, separately for each of the 335 full-support representatives.

\begin{proposition}[Exhaustive witness criterion]\label{prop:coverage}
For each canonical full-support premise triple $P$, every conclusion $c\in\A_5$ either belongs to the sound closure $C(P)$ or has a positive-definite Gaussian counterexample satisfying $P$.  Consequently, every regular-Gaussian-valid implication $P\Rightarrow c$ has $c\in C(P)$ and is a semantic consequence of $\Tfour$.
\end{proposition}

\begin{proof}
The rational certificates, accepted by the criterion below, give Gaussian roots satisfying $P$ and violating every atom in their cover sets.  The verified coverage equality therefore supplies a counterexample for each $c\notin C(P)$.  A valid conclusion must instead lie in $C(P)$, whose construction is sound for $\Tfour$.  Gaussian validity and $\Tfour$ consequence are invariant under $G$, so the same conclusion holds throughout each premise orbit.
\end{proof}

\subsection{Gaussian witness systems}

A correlation matrix on five labels has ten off-diagonal coordinates, ordered by pairs
\[
 (12),(13),(14),(15),(23),(24),(25),(34),(35),(45).
\]
For a premise triple $P=(p_1,p_2,p_3)$, each certificate designates three coordinates $y=(y_1,y_2,y_3)$ as unknowns and fixes the remaining seven at rational values.  Define
\[
 F(y)=\bigl(\AP_{p_1}(R(y)),\AP_{p_2}(R(y)),\AP_{p_3}(R(y))\bigr).
\]
The certificate consists of a rational box $X=X_1\times X_2\times X_3$, the fixed rational coordinates, and a nonsingular rational preconditioner $C$.  All polynomial and Jacobian interval evaluations use exact rational arithmetic.

For midpoint $x\in X$ and an interval enclosure $F'(X)$ of the Jacobian, define the Krawczyk operator~\cite{Krawczyk1969,Rump2010}
\begin{equation}\label{eq:krawczyk}
 K(x,X)=x-CF(x)+\bigl(I-CF'(X)\bigr)(X-x).
\end{equation}
The strict componentwise inclusion
\begin{equation}\label{eq:strictK}
 K(x,X)\subset\operatorname{int}(X)
\end{equation}
with nonsingular $C$ satisfies the standard Krawczyk existence-and-uniqueness criterion~\cite{Krawczyk1969,Rump2010}: $F$ has exactly one real zero in $X$.  Existence suffices for the classification; the criterion also proves uniqueness in the box.

At that root, positive definiteness follows from Sylvester's criterion.  The determinant of every leading principal block of the full interval matrix has a strictly positive lower endpoint.  For each conclusion atom listed in the certificate, interval evaluation of its almost-principal minor excludes zero.  The root is therefore a regular-Gaussian counterexample to every covered conclusion.

A single witness need not refute every conclusion outside the closure.  One certificate suffices for 306 canonical full-support representatives; for each of the remaining 29, the union of two certificates' cover sets equals $\A_5\setminus C(P)$.

\subsection{Worked certificate}

Certificate 272 (numbered from zero) treats the canonical premises
\[
 [12\mid3],\qquad[13\mid4],\qquad[14\mid5]
\]
and also supplies the premise-minimality witnesses in Corollary~\ref{cor:minimal}.  Its unknown coordinates are $r_{12},r_{13},r_{14}$.  The other seven correlations are fixed rational numbers.  Table~\ref{tab:worked} gives outward decimal enclosures; the supplementary data contain the exact rational values.

\begin{table}[ht]
\centering
\caption{Rational certificate 272.  Coordinate intervals are outward decimal enclosures of the exact boxes; the stated margins are computed from the original rational boxes.}
\label{tab:worked}
\begin{tabular}{lll}
\toprule
Quantity & Certified interval or lower bound & Role\\
\midrule
$r_{12}$ & $[0.000286218259399,\ 0.000286418259400]$ & $X_1\perp X_2$ is false\\
$r_{13}$ & $[-0.00536186144355,\ -0.00536166144354]$ & unknown root coordinate\\
$r_{14}$ & $[-0.0329011048000,\ -0.0329009047999]$ & unknown root coordinate\\
Krawczyk interior margin & $\ge 9.99\times10^{-8}$ in each coordinate & strict inclusion\\
$\det R_{1:2,1:2}$ & $>0.9999999179$ & Sylvester positivity\\
$\det R_{1:3,1:3}$ & $>0.9971197683$ & Sylvester positivity\\
$\det R_{1:4,1:4}$ & $>0.9462868919$ & Sylvester positivity\\
$\det R$ & $>0.7549461182$ & full positive definiteness\\
\bottomrule
\end{tabular}
\end{table}

The fixed correlations, in pair order $(15),(23),(24),(25),(34),(35),(45)$, are approximately
\[
 (-0.231821,-0.053400,0.144036,-0.179468,0.162966,-0.256434,0.141924).
\]
The supplementary data give the rational box, all seven fractions, the $3\times3$ rational matrix $C$, and certified nonvanishing for 77 atoms outside the closure.

\FloatBarrier
\subsection{Proper-support extension lemma}

The finite certificate enumeration is needed only when the premises collectively use all five labels.

\begin{lemma}[Generic positive-definite extension]\label{lem:extension}
Let $P$ be any finite set of elementary CI atoms supported on a proper subset $S\subsetneq N$.  If a proposed conclusion $c$ uses a label outside $S$, then there is a positive-definite rational covariance matrix satisfying every atom in $P$ and violating $c$.  One matrix may be chosen to violate any prescribed finite collection of such conclusions simultaneously.
\end{lemma}

\begin{proof}
Set the principal covariance block on $S$ equal to the identity.  Then every atom supported on $S$, in particular every premise, holds.  Extend to all labels with
\[
 \Sigma=I+E,
\]
where the $S\times S$ block of $E$ is zero and only covariances involving labels outside $S$ may vary.  For sufficiently small $E$, for example $\|E\|_2<1$, all eigenvalues of $I+E$ are positive.  The old principal block is unchanged, so all premises remain true.

The almost-principal minor $g_c(E)$ for a conclusion involving an outside label is a polynomial in the new entries, and it is not identically zero.  If an endpoint of $c$ lies outside $S$, set only the direct covariance between the two endpoints to a nonzero parameter $t$ and leave conditioned labels independent.  Then $g_c(E)=\pm t$.  If both endpoints lie in $S$ and an outside label $x$ appears only in the conditioning set, correlate $x$ with both endpoints by parameters $u,v$ and leave all other new entries zero.  The Schur-complement cross term is $-uv$, so $g_c(E)$ is nonzero for suitable $u,v$.

Thus each unwanted conclusion defines a proper algebraic subset of an open neighborhood of $E=0$.  A finite union of proper real algebraic subsets cannot fill that neighborhood.  Rational points are dense, so a sufficiently small rational $E$ avoids all of them while preserving positive definiteness and every premise.
\end{proof}

If both premises and conclusion lie in $S$, any regular Gaussian on $S$ extends independently to five labels.  Conversely, embed $S$ in a four-label marginal face, adding unused labels when $|S|<4$.  A valid implication on $S$ holds in every regular-Gaussian structure in the complete four-variable catalogue and hence in that face of every $\Tfour$ model.  It is therefore $\Tfour$-local.  Lemma~\ref{lem:extension} shows that $P\Rightarrow c$ is Gaussian-invalid whenever $c$ uses a label outside $S$.  Thus every rule at any of the 104 canonical proper-support representatives is invalid or $\Tfour$-local.  Orbit invariance extends this conclusion to every member of the orbit, including members with full support.

\subsection{Zero, one, and two antecedents}

Boege's two-antecedental completeness theorem states that every inference form with at most two antecedents that is valid for regular Gaussians on any finite ground set follows from the gaussoid axioms~\cite{Boege2022,BoegeCorrection2022}.  Horn implications are a special case.  This includes zero antecedents; directly, no elementary atom is universally valid, since a correlation matrix with only $r_{ij}=r_{ji}=1/2$ nonzero off the diagonal violates $[ij\mid K]$ for every admissible $K$.  As observed after Definition~\ref{def:t4}, each five-variable gaussoid-axiom instance lies in a marginal or conditional four-face.  Hence every valid rule with at most two antecedents is a semantic consequence of $\Tfour$.

Proposition~\ref{prop:coverage} handles the full-support representatives; Lemma~\ref{lem:extension} and the four-face argument handle the proper-support representatives.  Together with the low-antecedent theorem and symmetry reduction, these prove Proposition~\ref{prop:lower}, and with Proposition~\ref{prop:upper} prove Theorem~\ref{thm:main}.

\section{Verification and computational methods}\label{sec:verification}

The finite enumerations and rational interval inequalities in Table~\ref{tab:architecture} are checked in exact arithmetic.  Numerical search is used only to locate candidate boxes.

\begin{table}[ht]
\centering
\caption{Claim-to-certificate map.}
\label{tab:architecture}
\small
\begin{tabular}{P{3.0cm}P{5.4cm}P{3.8cm}}
\toprule
Claim & Proof object & Arithmetic domain\\
\midrule
$|\A_5|=80$, symmetry order 240 & atom enumeration and permutation/duality maps & finite integer combinatorics\\
629 retained four-variable models & exhaustive axiom-model enumeration and five additional implication schemes & finite bit logic\\
Upper-bound validity & telescoping identity~\eqref{eq:telescoping} & integer polynomials\\
Upper-bound nonlocality & 80-bit countermodel and ten face masks & finite bit logic\\
Three-antecedent exhaustion & 439 premise orbits; 105 mixed-support; 21,357 rule orbits in the certified stratum & finite group action\\
Gaussian counterexamples & 364 Krawczyk certificates and their conclusion cover sets & exact rational intervals\\
Positive definiteness & five leading-principal determinant enclosures (the $1\times1$ block is the unit diagonal) & exact rational intervals\\
Conclusion nonvanishing & almost-principal determinant enclosures & exact rational intervals\\
Premise minimality & four deletion mappings to certificate 272 & finite symmetry plus intervals\\
\bottomrule
\end{tabular}
\end{table}

For each witness, the proof checks a nonsingular rational preconditioner, strict Krawczyk inclusion, positive leading-principal determinant enclosures, and nonzero conclusion-minor enclosures.  Coverage is checked against the complement of each computed closure.  Exhaustive enumeration of the CNF models gives 679 gaussoids; the five additional implication schemes retain 629, identified with the regular-Gaussian structures by the classification theorem in~\cite{LM2007}.

\subsection{AI-assisted research and verification}

OpenAI GPT-5.6 Sol, GPT-6 Astra, GPT-6.1 Sol, and GPT-6.1 Terra, and DeepSeek Flash and DeepSeek V4 Pro assisted with literature exploration, conjecture and proof development, computational implementation and checking, and manuscript drafting and revision.  DeepSeek Flash was used as an implementation subagent; GPT-6.1 Terra and GPT-6.1 Sol were also used in the preparation of this revision.  AI-generated suggestions are not mathematical evidence; the computational claims rely on the exact-arithmetic verifier and certificate data described above.

\section{Related work}\label{sec:prior}

Mat\'u\v{s} formulated Gaussian conditional independence through polynomial relations among covariance minors~\cite{Matus2005}.  Ln\v{e}ni\v{c}ka and Mat\'u\v{s} subsequently classified the four-variable regular-Gaussian structures and gave the five additional implication schemes used here to reconstruct $\Tfour^{(4)}$~\cite{LM2007}.

The general shared-endpoint implication family \eqref{eq:family} is due to \v{S}ime\v{c}ek~\cite{Simecek2006}; its $k=4$ instance is \eqref{eq:cycle5}.  The general family is therefore prior work.  The present argument adds the $\Tfour$ countermodel, premise minimality, and the exhaustive three-antecedent lower bound for its five-variable member.

Sullivant established non-finite characterization by a distinct cyclic family~\cite{Sullivant2009}.  Its premise endpoints form a cycle, whereas the premises in \eqref{eq:family} share one endpoint; their endpoint graphs have different degree sequences and are not related by relabelling.  Sullivant's proof uses binomial primary decomposition and excludes components from the positive-definite cone by Hadamard-product arguments, while \eqref{eq:telescoping} gives a direct proof for the family needed here.

Boege proved two-antecedental completeness of the gaussoid axioms, with a later correction, and introduced selfadhesivity as a stronger necessary approximation to Gaussian realizability~\cite{Boege2022,BoegeCorrection2022,BoegeSA2023}.  Related work on the geometry and construction of gaussoids further distinguishes abstract gaussoid structure from positive-definite Gaussian realizability~\cite{Geometry2019,Construction2020}.

\section{Discussion and open problems}\label{sec:discussion}

Theorem~\ref{thm:main} determines the least antecedent count but not the full set of four-antecedent rules beyond $\Tfour$.  Several related problems remain.
\begin{enumerate}
 \item A classification of premise-minimal four-antecedent rules remains open.
 \item A structural proof of the three-antecedent lower bound could replace or compress the 364-certificate exhaustion.
 \item For $n>5$, one can define the threshold \(a_n\) relative to the complete $(n-1)$-local face theory.  \v{S}ime\v{c}ek's family supplies possible upper bounds, but nonlocality and lower bounds must be proved separately.
\end{enumerate}

The classification in this paper is limited to elementary CI atoms, singleton-conclusion Horn implications, regular positive-definite Gaussians, five labels, and at most three antecedents in the lower-bound search.  It does not cover singular Gaussians, disjunctive conclusions, or general semigraphoids.

\paragraph{Reproducibility.}
The certificate data and exact-arithmetic verifier are available in software release \href{https://github.com/Maar10Herr/Gaussian-CI-a5-Equals-4-Public/tree/v0.3.1}{v0.3.1}.  From the repository root, \code{make verify-core} checks the finite enumeration, polynomial identity, countermodel, and rational interval certificates.

\href{https://github.com/Maar10Herr/Gaussian-CI-a5-Equals-4-Public/tree/af3848c44caabbdba891dab7efbdf9c7ee81d4c0/formalization}{Supplementary Lean proofs} establish the positive-definite covariance-matrix implication~\eqref{eq:cycle5}, a four-local countermodel, and premise minimality.

\paragraph{Acknowledgment.}
I thank Tobias Boege for helpful correspondence on Gaussian conditional independence and the relevant literature.

\appendix
\section{Canonical atom order and countermodel}\label{app:atoms}

The 80 atoms are ordered deterministically.  Iterate unordered endpoint pairs $(i,j)$ lexicographically.  For each pair, list the remaining labels increasingly as $(u_0,u_1,u_2)$ and iterate masks $m=0,\ldots,7$; bit $t$ of $m$ includes $u_t$ in $K$.  Thus indices $0,1,2,3$ begin
\[
 [12],\ [12\mid3],\ [12\mid4],\ [12\mid34],
\]
and the same mask convention continues through pair $(4,5)$.  This convention uniquely determines every index.

The nonlocality countermodel is
\begin{center}
\ttfamily\small
01000000001000000000101001000000000000000000000000001001001011000000000000000000
\end{center}
where one means membership in the abstract CI structure.

\section{Construction of the four-variable catalogue}\label{app:t4}

Translating every elementary gaussoid-axiom instance on the 24 four-variable atoms into propositional clauses and enumerating all satisfying assignments gives 679 models.  Each model is tested, under every permutation of four labels, against the five additional implication schemes numbered (17)--(21) by Ln\v{e}ni\v{c}ka and Mat\'u\v{s}~\cite{LM2007}.  Exactly 50 models violate at least one scheme, leaving 629.  By the classification theorem in~\cite{LM2007}, these are precisely the regular-Gaussian structures.

\section{Certificate data}\label{app:cert}

A certificate specifies three premise indices, seven fixed rational correlations, the three remaining coordinate indices, a rational box for those coordinates, a nonsingular rational $3\times3$ preconditioner, and a set of covered conclusions.  The fixed and unknown coordinates partition the ten off-diagonal entries.  Determinant and derivative enclosures are recomputed by signed permutation expansion.  The complete fractions are recorded in \code{m3\_interval\_certificates.json} in the repository.

For certificate 272 the first rational unknown box is
\[
 \left[\frac{1036593930957}{3621690430000000},
       \frac{1037318269043}{3621690430000000}\right],
\]
which excludes zero and therefore refutes the marginal independence $X_1\perp X_2$.

\bibliographystyle{plain}
\bibliography{references}

\end{document}
